\documentclass[12 pt, letterpaper]{hmcpset}
\usepackage{hw-preamble}
\usepackage{pdflscape}
\setlist[enumerate, 1]{label = (\roman*)}

\name{Jayden Thadani}
\class{MATH 145 --- Topics in Geometry and Topology}
\assignment{Homework 7}
\duedate{11th March 2025}

\begin{document}

\begin{problem}[34.]
	Let $G = (V, E)$ be a finite graph.
	\begin{enumerate}
		\item Show that
			\[
				\sum_{x \in V} \lvert N(x) \rvert \Delta f(x) = 0
			\]
			for any scalar field $f \in \mathrm{C}^{0}(G)$.

		\item Suppose the graph is connected. Let $f \in \mathrm{C}^{0}(G)$ be a non-constant scalar field. Show that there exist points $x, y \in V$ such that $\Delta f(x) > 0$ and $\Delta f(y) < 0$.

		\item Show that (ii) fails if the graph isn't connected.
	\end{enumerate}
\end{problem}

\begin{solution}
	\begin{enumerate}
		\item Since
			 \[
				\Delta f(x) = \sum_{y \in N(x)} \frac{f(x) - f(y)}{N(x)}
			\]
			we can apply the handshake lemma to get
			\begin{align*}
				\sum_{x \in V} \lvert N(x) \rvert \Delta f(x) & = \sum_{x \in V} \sum_{y \in N(x)} f(x) - f(y) \\
									      & = \sum_{x \in V} \operatorname{deg} x \, f(x) - \sum_{y \in V} \deg y \, f(y) \\
									      & = 0.
			\end{align*}
			That is, the second equality follows because every vertex is counted as the neighbour of other vertices as many times as its degree.

		\item We showed in class that if $\Delta f = 0$ everywhere on a connected graph, then $f$ is constant. Thus, if $f$ is non-constant, there must be some point $x_{0}$ such that $\Delta f(x_{0}) > 0$ or some point $y_{0}$ such that $\Delta f(y_{0}) < 0$. Suppose without loss of generality that we have $x$ with $\Delta f(x_{0}) > 0$.

			For the sum
			\[
				\sum_{x \in V} \lvert N(x) \rvert \Delta f(x)
			\]
			to be zero, there must be some negative term $\lvert N(y_{0}) \rvert \Delta f(y_{0}) < 0$. For each $x \in V$ $\lvert N(x) \rvert \ge 0$ since it is the cardinality of the set. Thus, $\Delta f(y_{0}) < 0$.

		\item Any non-constant, locally constant function on a disconnected graph is harmonic everywhere. For example, on the disjoint union of two two-point graph $G(V) = \{ 1, 2 \} \sqcup \{ 3, 4 \}$ with only edges $G(E) = \{ 12, 34 \}$, we can define a function
			\[
				f(x) = \begin{cases}
					1 & x \le 2 \\
					0 & 3 \le x.
				\end{cases}
			\]
			At each vertex $x$, the value of $f$ minus the average value of $f$ on $N(x)$ is $0$ since $f$ is constant on $N(x)$. However, the function isn't globally constant. Essentially, $\Delta f = 0$ forces $f$ to be locally constant, which doesn't force $f$ to be globally constant on non-connected graph.x
	\end{enumerate}
\end{solution}

\pagebreak

\begin{problem}[35.]
	Consider the Dirichlet problem for the following graph and boundary values:
	\begin{figure}[H]
		\centering
		\includegraphics[width = 0.3 \textwidth]{figures/P35 graph.pdf}
	\end{figure}
	\begin{enumerate}
		\item Express the Dirichlet problem as a system of five equations in five unknowns.
		
		\item Solve the Dirichlet problem.
	\end{enumerate}
\end{problem}

\begin{solution}
	We label the interior vertices and then the boundary vertices clockwise from bottom-left as in the image below.

	\begin{figure}[H]
		\centering
		\includegraphics[width = 0.3 \textwidth]{figures/P35 labels}
		\caption{Labels for the vertices of the graph in blue}
	\end{figure}

	Thus, we have the following matrix representing the linear map $D : f \mapsto (\Delta f_{\mid I}, f_{\mid B})$ ---
	\[
		\begin{pmatrix}
			1 & -1/2 & 0 & 0 & 0 & 0 & 0 & 0 & -1/2 \\
			-1/4 & 1 & -1/4 & 0 & -1/4 & -1/4 & 0 & 0 & 0 \\
			0 & -1/2 & 1 & -1/2 & 0 & 0 & 0 & 0 & 0 \\
			0 & 0 & -1/3 & 1 & -1/3 & 0 & -1/3 & 0 & 0 \\
			0 & -1/4 & 0 & -1/4 & 1 & 0 & 0 & -1/4 & -1/4 \\
			0 & 0 & 0 & 0 & 0 & 1 & 0 & 0 & 0 \\
			0 & 0 & 0 & 0 & 0 & 0 & 1 & 0 & 0 \\
			0 & 0 & 0 & 0 & 0 & 0 & 0 & 1 & 0 \\
			0 & 0 & 0 & 0 & 0 & 0 & 0 & 0 & 1 \\
		\end{pmatrix}.
	\]
	We want to solve
	\[
		\begin{pmatrix}
			1 & -1/2 & 0 & 0 & 0 & 0 & 0 & 0 & -1/2 \\
			-1/4 & 1 & -1/4 & 0 & -1/4 & -1/4 & 0 & 0 & 0 \\
			0 & -1/2 & 1 & -1/2 & 0 & 0 & 0 & 0 & 0 \\
			0 & 0 & -1/3 & 1 & -1/3 & 0 & -1/3 & 0 & 0 \\
			0 & -1/4 & 0 & -1/4 & 1 & 0 & 0 & -1/4 & -1/4 \\
			0 & 0 & 0 & 0 & 0 & 1 & 0 & 0 & 0 \\
			0 & 0 & 0 & 0 & 0 & 0 & 1 & 0 & 0 \\
			0 & 0 & 0 & 0 & 0 & 0 & 0 & 1 & 0 \\
			0 & 0 & 0 & 0 & 0 & 0 & 0 & 0 & 1 \\
		\end{pmatrix} \begin{pmatrix}
			f(v_{1}) \\
			\vdots \\
			f(v_{5}) \\
			f(v_{6}) \\
			\vdots \\
			f(v_{9})
		\end{pmatrix} = \begin{pmatrix}
			0 \\
			\vdots \\
			0 \\
			0 \\
			3 \\
			1 \\
			2
		\end{pmatrix}.
	\]
	With the knowledge that $f(v_{6}), \dots, f(v_{9})$ must be $0, 3, 1, 2$ respectively, we can reduce this to the linear system on $\mathbb{R}^{5}$ ---
	\begin{align*}
		\begin{pmatrix}
			1 & -1/2 & 0 & 0 & 0 \\
			-1/4 & 1 & -1/4 & 0 & -1/4 \\
			0 & -1/2 & 1 & -1/2 & 0 \\
			0 & 0 & -1/3 & 1 & -1/3 \\
			0 & -1/4 & 0 & -1/4 & 1 \\
		\end{pmatrix} \begin{pmatrix}
			f(v_{1}) \\
			\vdots \\
			f(v_{5})
			\end{pmatrix} & = - \begin{pmatrix}
			0 & 0 & 0 & -1/2 \\
			-1/4 & 0 & 0 & 0 \\
			0 & 0 & 0 & 0 \\
			0 & -1/3 & 0 & 0 \\
			0 & 0 & -1/4 & -1/4 \\
		\end{pmatrix} \begin{pmatrix}
			0 \\
			3 \\
			1 \\
			2
		\end{pmatrix} \\
		& = \begin{pmatrix}
			1 \\
			0 \\
			0 \\
			1 \\
			3/4
		\end{pmatrix}
	\end{align*}

	Solving using an online matrix solver gives
	\[
		\begin{pmatrix}
			f(v_{1}) \\
			f(v_{2}) \\
			f(v_{3}) \\
			f(v_{4}) \\
			f(v_{5})
		\end{pmatrix} = \begin{pmatrix}
			8/5 \\
			6/5 \\
			49/30 \\
			31/15 \\
			47/30
		\end{pmatrix}.
	\]
\end{solution}

\pagebreak

\begin{problem}[36.]
	Let $G = (V, E, B)$ be a finite graph-with-boundary with $\mathrm{b}_{0}(G, B) = 0$ and let $h : B \to \mathbb{R}$ be a system of boundary values. This gives rise to a Dirichlet problem. Suppose the Dirichlet problem has a \emph{symmetry}. This is a bijection $\alpha : V \to V$ with the following properties:
	\begin{itemize}
		\item It preserves the graph structure, meaning that $\alpha(x) \alpha(y) \in E$ if and only if $xy \in E$.
		\item It preserves the boundary, meaning that $\alpha(x) \in B$ if and only if $x \in B$.
		\item It preserves the boundary values, meaning that $g(\alpha(x)) = g(x)$ for all $x \in B$.
	\end{itemize}
	Let $f : V \to \mathbb{R}$ be a solution to the Dirichlet problem. Give a two-sentence proof that the symmetry preserves $f$, meaning that $f(\alpha(x)) = f(x)$ for all $x \in V$.
\end{problem}

\begin{solution}
	Since $\alpha$ preserves edges, it preserves the Laplacian of any function with $\alpha$-symmetry, specifically including $g$ ---
	\[
		\Delta g(\alpha(x)) = \sum_{\alpha(y) \in N(\alpha(x))} \frac{g(\alpha(x)) - g(\alpha(y))}{\lvert N(\alpha(x)) \rvert} = \sum_{y \in N(x)} \frac{g(x) - g(y)}{\lvert N(x) \rvert}.
	\]
	But then, if $f' = f \circ \alpha \in \mathrm{C}_{0}(G, B)$ has $\Delta(f') = 0$ on the interior and $f' = g \circ \alpha$ on the boundary, $f$ must be the unique solution to the Dirichlet problem with $\Delta f = 0$ on the interior and $f = g$ on the boundary.
\end{solution}

\pagebreak

\begin{landscape}

\begin{problem}[37.]
	Consider the Dirichlet problem for the following graph and boundary values:
	\begin{figure}[H]
		\centering
		\includegraphics[width = 0.3 \textwidth]{figures/P37 graph.pdf}
	\end{figure}
	\begin{enumerate}
		\item Express the Dirichlet problem as a linear system of eight equations in eight unknowns.
		\item Express the Dirichlet problem as a linear system of two equations in two unknowns.

		\item Solve the Dirichlet problem.
	\end{enumerate}
\end{problem}

\begin{solution}
	We label the vertices anticlockwise from middle of the right, first the interior vertices, then the boundary vertices, with the central vertex labelled last.
	\begin{enumerate}
		\item The Dirichlet problem corresponds to solving the following linear system (using the techniques from homework 8 problem 41)
			\setcounter{MaxMatrixCols}{20}
			\begin{align*}
				& \begin{pmatrix}
   1 & -\frac{1}{4} &    0 &    0 &    0 &    0 &    0 & -\frac{1}{4} \\
-\frac{1}{4} &    1 & -\frac{1}{4} &    0 &    0 &    0 &    0 &    0 \\
   0 & -\frac{1}{4} &    1 & -\frac{1}{4} &    0 &    0 &    0 &    0 \\
   0 &    0 & -\frac{1}{4} &    1 & -\frac{1}{4} &    0 &    0 &    0 \\
   0 &    0 &    0 & -\frac{1}{4} &    1 & -\frac{1}{4} &    0 &    0 \\
   0 &    0 &    0 &    0 & -\frac{1}{4} &    1 & -\frac{1}{4} &    0 \\
   0 &    0 &    0 &    0 &    0 & -\frac{1}{4} &    1 & -\frac{1}{4} \\
-\frac{1}{4} &    0 &    0 &    0 &    0 &    0 & -\frac{1}{4} &    1
\end{pmatrix} \begin{pmatrix}
	f(v_{1}) \\
	\vdots \\
	f(v_{8})
\end{pmatrix} \\ & = \begin{pmatrix}
-\frac{1}{4} &    0 & 0 &    0 &    0 &    0 & 0 &    0 &    0 &    0 & 0 &    0 &    0 &    0 & 0 &    0 & -\frac{1}{4} \\
   0 & -\frac{1}{4} & 0 & -\frac{1}{4} &    0 &    0 & 0 &    0 &    0 &    0 & 0 &    0 &    0 &    0 & 0 &    0 &    0 \\
   0 &    0 & 0 &    0 & -\frac{1}{4} &    0 & 0 &    0 &    0 &    0 & 0 &    0 &    0 &    0 & 0 &    0 & -\frac{1}{4} \\
   0 &    0 & 0 &    0 &    0 & -\frac{1}{4} & 0 & -\frac{1}{4} &    0 &    0 & 0 &    0 &    0 &    0 & 0 &    0 &    0 \\
   0 &    0 & 0 &    0 &    0 &    0 & 0 &    0 & -\frac{1}{4} &    0 & 0 &    0 &    0 &    0 & 0 &    0 & -\frac{1}{4} \\
   0 &    0 & 0 &    0 &    0 &    0 & 0 &    0 &    0 & -\frac{1}{4} & 0 & -\frac{1}{4} &    0 &    0 & 0 &    0 &    0 \\
   0 &    0 & 0 &    0 &    0 &    0 & 0 &    0 &    0 &    0 & 0 &    0 & -\frac{1}{4} &    0 & 0 &    0 & -\frac{1}{4} \\
   0 &    0 & 0 &    0 &    0 &    0 & 0 &    0 &    0 &    0 & 0 &    0 &    0 & -\frac{1}{4} & 0 & -\frac{1}{4} &    0 \\
\end{pmatrix} \begin{pmatrix}
	0 \\
	\vdots \\
	0 \\
	1
\end{pmatrix}
			\end{align*}
		\item Notice that the graph has fourfold symmetry --- the natural action of the cyclic group $Z_{4}$ on the graph is symmetric. Thus, we only have to solve the Dirichlet problem in one corner of the graph. We can write the same linear map $D : f \mapsto (\Delta f, f)$ as previously now as
			\[
				\begin{pmatrix}
					1 & -1/2 \\
					-1/2 & 1
				\end{pmatrix} \begin{pmatrix}
					f(v_{1}) \\
					f(v_{2})
				\end{pmatrix} = \begin{pmatrix}
					-1/4 & 0 & -1/4 \\
					0 & -1/2 & 0
				\end{pmatrix} \begin{pmatrix}
					0 \\
					0 \\
					1
				\end{pmatrix}
			\]
		\item 
			\[
				\boxed{
					f(v_{n}) = \begin{cases}
						1/6 & n = 2k \\
						1/3 & n = 2k + 1
					\end{cases}
				}
			\]

			The matrix equation in (ii) reduces to
			\[
				\begin{pmatrix}
					1 & -1/2 \\
					-1/2 & 1
				\end{pmatrix} \begin{pmatrix}
					f(v_{1}) \\
					f(v_{2})
				\end{pmatrix} = \begin{pmatrix}
					-1/4 \\
					0
				\end{pmatrix}
			\]
			which has the unique solution $f(v_{1}) = 1/3$ and $f(v_{2}) = 1/6$. Extending the results to the whole graph by symmetry gives the result above.
\end{enumerate}
\end{solution}	
\end{landscape}

\pagebreak

\begin{problem}[38.]
	Consider the `infinite plane grid' graph whose vertices are the points $(m, n)$ where $m, n \in \mathbb{Z}$, and where the two vertices are connected by an edge if and only if the two points are distance $1$ apart.
	\begin{enumerate}
		\item We define six scalar fields
			\begin{align*}
				f(m, n) & = 1 & g_{1}(m, n) & = m & g_{2}(m, n) & = n \\
				h_{1}(m, n) & = m^{2} & h_{2}(m, n) & = mn & h_{3}(m, n) & = n^{2}.
			\end{align*}

			Calculate $\Delta f$, $\Delta g_{1}$, $\Delta g_{2}$, $\Delta h_{1}$, $\Delta h_{2}$, and $\Delta h_{3}$.

		\item Show that the space of harmonic scalar fields of degree at most $2$ is five-dimensional and give a basis for that space.
	\end{enumerate}
\end{problem}

\begin{solution}
	Note that the neighbours of each point $(m, n) \in \mathbb{Z}^{2}$ are
	\[
		N(m, n) = \{ (m + 1, n), (m - 1, n), (m, n + 1), (m, n - 1) \}.
	\]
	Also, we use the formula
	\[
		\Delta f(m, n) = \frac{\sum_{(x, y) \in N(m, n)} f(m, n) - f(x, y)}{\lvert N(m, n) \rvert}.
	\]
	\begin{enumerate}
		\item Since $f$ is constant, $\Delta f = 0$ everywhere.

			The two linear functions are both harmonic ---
			\begin{align*}
				\Delta g_{1}(m, n) & = \frac{(m - (m - 1)) + (m - (m + 1)) + (m - m) + (m - m)}{4} \\
						   & = \frac{1 - 1 + 0 + 0}{4} \\
						   & = 0
			\end{align*}
			and
			\begin{align*}
				\Delta g_{2}(m, n) & = \frac{(n - (n - 1)) + (n - (n + 1)) + (n - n) + (n - n)}{4} \\
						   & = \frac{1 - 1 + 0 + 0}{4} \\
						   & = 0.
			\end{align*}
			The product $h_{2}$ is harmonic ---
			\begin{align*}
				\Delta h_{2}(m, n) & = - \frac{((m + 1) n - mn) + ((m - 1)n - mn) + (m(n + 1) - mn) + (m(n - 1) - mn)}{4} \\
						   & = - \frac{n - n + m - m}{4} \\
						   & = 0.
			\end{align*}
			However, the square functions are not ---
			\begin{align*}
				\Delta h_{1}(m, n) & = - \frac{((m + 1)^{2} - m^{2}) + ((m - 1)^{2} - m^{2}) + (m^{2} - m^{2}) + (m^{2} - m^{2})}{4} \\
					     & = - \frac{(2m + 1) - (2m - 1)}{4} \\
					     & = - \frac{1}{2}.
			\end{align*}
			and similarly, $\Delta h_{3}(m, n) = -1/2$.

		\item The space of all scalar fields of degree at most $2$ is spanned by the basis vectors $f, g_{1}, g_{2}, h_{1}, h_{2}, h_{3}$. Clearly $\Delta h_{1}(0, 0) \neq 0$ so $h_{1}$ is not harmonic. Thus, the space of harmonic scalar fields is at most $(6 - 1)$-dimensional. We will demonstrate a list of five linearly independent vectors, and thus, show that the space is exactly $5$-dimensional. 


			The list $f, g_{1}, g_{2}, h_{2}$ is linearly independent (since they're part of the basis of the space of functions we're considering) and each function is harmonic (by our previous calculations). Further, since $\Delta h_{1} = \Delta h_{3}$, and $\Delta$ is linear, the function $h_{1} - h_{3}$ is harmonic. This is also linearly independent of the other $4$ functions since the linear dependence of $h_{1} - h_{3}$ on $f, g_{1}, g_{2}, h_{2}$ would mean that the basis $f, g_{1}, g_{2}, h_{1}, h_{2}, h_{3}$ is not linearly independent. Thus, $f, g_{1}, g_{2}, h_{1} - h_{3}, h_{2}$ is a list of five linearly independent harmonic functions on $\mathbb{Z}^{2}$. Since we've shown already that the space of harmonic functions is at most $5$-dimensional, this is a basis of the space of harmonic functions.
	\end{enumerate}
\end{solution}

\pagebreak

\begin{problem}[38.]
	Analogous to the previous question, consider the space of smooth function $\mathbb{R}^{2} \to \mathbb{R}$, and the Laplacian differential operator
	\[
		\Delta = - \left ( \frac{\partial^{2}}{\partial x^{2}} + \frac{\partial^{2}}{\partial y^{2}} \right ).
	\]
	\begin{enumerate}
		\item Show that the space of harmonic functions of degree at most $2$ is five-dimensional and give a basis for that space.

		\item Show that the $(m + 1)$-dimensional space of homogeneous polynomials of degree $m$,
			\[
				\mathcal{HP}_{m} = \{ a_{0} x^{m} + a_{1}^{m - 1} x^{m - 1} y + \dots + a_{m} y^{m} \mid a_{0}, a_{1}, \dots, a_{m} \in \mathbb{R} \}
			\]
			contains an at least $2$-dimensional subspace of functions that are harmonic.
	\end{enumerate}
\end{problem}

\begin{enumerate}
	\item The space of harmonic polynomials of degree at most $2$ (which we write as $\ker \Delta$) cannot be a full subspace of $\mathcal{P}_{2}$ (the space of all polynomials in $x$ and $y$ of degree at most $2$) since there are non-harmonic polynomials in $\mathcal{P}_{2}$
		\[
			\Delta(-x^{2}) = -1 \neq 0.
		\]
		Thus, $\ker \Delta$ is at most $5$-dimensional. We will show that it contains $5$ linearly independent harmonic polynomials, and thus, is exactly $5$-dimensional with those polynomials as a basis.

		The polynomials $1$, $x$, $y$ are clearly harmonic since they are sub-quadratic and vanish on application of any second-order differential operators. They are also linearly independent since they form a basis of the polynomials of degree at most $1$. Further, by computing
		\[
			- \Delta(xy) =  \frac{\partial^{2} (xy)}{\partial x^{2}} + \frac{\partial^{2} (xy)}{\partial y^{2}}= 0 + 0 = 0
		\]
		we see that $xy$ is also harmonic. The list $1, x, y, xy$ is still linearly independent since it is a sublist of the basis of $\mathcal{P}_{2}$. Finally,
		\[
			- \Delta(x^{2} + y^{2}) = \frac{\partial^{2} x^{2}}{\partial x^{2}} - \frac{\partial^{2} y^{2}}{\partial y^{2}} = 1 - 1 = 0.
		\]
		so $x^{2} - y^{2}$ is harmonic. The list $1, x, y, xy, x^{2} - y^{2}$ is linearly independent since the list $1, x, y, xy, x^{2}, y^{2}$ is the basis of $\mathcal{P}_{2}$.

	\item Since $\Delta$ is a differential operator of order $2$ and homogeneous (each linear term is of the same), it drops the degree of any polynomial by $2$. Thus, the image of $\mathcal{HP}_{m}$ under $\Delta$ is contained in $\mathcal{HP}_{m - 2}$. This is $2$ dimensions smaller than $\mathcal{HP}_{m}$, so by the rank-nullity theorem, $\Delta$ has a kernel of dimension at least $2$.
\end{enumerate}

\end{document}
